% T2.3 / R2-3 — Section IV-E, "Cost inside an attention pipeline".
% Reviewer 2: "report latency and integration overhead inside a complete
% attention pipeline."
%
% The cycle-budget and MAC-ratio arguments below are exact and need no
% external data. The one number that must come from a cited paper is the
% area fraction against a published attention block -- marked <<TBD>>; take it
% from ITA [c8] rather than estimating.

\subsection{Cost Inside an Attention Pipeline}
\label{sec:integration}

A rotation datapath is only useful if it keeps up with the attention block it
feeds, and cheap enough that it does not dominate what it accelerates. Both
can be settled from first principles.

\subsubsection{Where the angle front-end amortises}
Each dimension pair has its own frequency, $\theta_i = m\Theta_i$, so every
lane requires its own angle: the front-end is not shared across the $d_h/2$
lanes of a head. What the angle does not depend on is the \emph{data}. The
digit string for a given $(m, i)$ is therefore identical in every attention
head, every layer, and every batch element, and can be computed once and
broadcast. For a LLaMA-2-7B-class model that is $32$ heads $\times$ $32$
layers $= 1024$ uses of each encoded angle, so the encoder and range-reduction
cost enters the per-rotation budget divided by $1024$. This is the property
that makes a feed-forward encoder affordable, and it is also why the
multiplier baseline of Table~\ref{tab:perf} is not a complete design: it
consumes $\sin$ and $\cos$ that something must still generate and store.

\subsubsection{Cycle budget}
Per head, layer and generated token, RoPE performs $d_h/2 = 64$ rotations for
the query and $64$ for the new key --- $128$ in total. Over the same interval
the attention block computes $QK^\top$ and the value product, $2 \cdot
N_{\text{ctx}} \cdot d_h$ multiply--accumulates. Table~\ref{tab:budget} gives
the resulting lane requirement at $N_{\text{ctx}} = 2048$.

\begin{table}[t]
\caption{RoPE lanes required to keep pace with attention, per head and layer,
at $N_{\mathrm{ctx}} = 2048$, $d_h = 128$. One pipelined lane retires one
rotation per cycle.}
\label{tab:budget}
\centering
\begin{tabular}{rrr}
\toprule
Attention MACs & Cycles per head-token & RoPE lanes needed \\
\midrule
64  & 4096 & 0.03 \\
128 & 2048 & 0.06 \\
256 & 1024 & 0.13 \\
512 & 512  & 0.25 \\
\bottomrule
\end{tabular}
\end{table}

Even a narrow $512$-MAC attention datapath leaves $512$ cycles in which to
retire $128$ rotations, so a \emph{single} pipelined lane is already four
times over-provisioned. The rotator is never the throughput bottleneck, which
is what makes the latency figures of Table~\ref{tab:perf} unimportant in
practice: a $4.48$\,ns pipeline traversal is hidden entirely behind an
attention step that takes hundreds of cycles.

\subsubsection{Share of the arithmetic}
The same comparison in multiply-equivalents: attention performs
$524\,288$ MACs per head, layer and token at $N_{\text{ctx}} = 2048$, against
$128$ rotations. Implemented with multipliers those rotations would be $512$
MAC-equivalents, or $0.098\%$ of the attention arithmetic; implemented as
shift-and-add they contribute no multipliers at all. Against a published
quantised attention accelerator~\cite{c8}, the proposed unit occupies
<<TBD>>\% of the block area.

The practical reading is that the choice of rotation datapath is not a
throughput decision and barely an area one --- it is an energy and
composability decision. At $463$\,nJ per token (Table~\ref{tab:perf}) against
$1191$\,nJ for the LUT alternative, the difference is worth having, and the
multiplier-free datapath removes a multiplier array from a block that an
attention accelerator would otherwise have to schedule around.
